\chapter{Repository Information}
\label{chap:repository}

All simulation pipelines, evaluation scripts, and structural configurations developed for this thesis are version-controlled and publicly accessible. This ensures full reproducibility of the 1.8 million dataset simulations discussed in the preceding chapters.

\section{Access Information}
\begin{itemize}
    \item \textbf{Repository URL:} \url{https://github.com/Sterrysx/ML_SyntheticDataResearch}
    \item \textbf{Primary Author:} Oriol Farr\'es (Supervised by Jordi Cort\'es)
    \item \textbf{Status:} Production-ready pipeline
    \item \textbf{Original Work Base:} Chen Xinnuo (2025)
\end{itemize}

\section{Repository Architecture and Documentation}
Rather than reproducing raw documentation, this section synthesizes the core structure and operational guides of the repository to provide a clear technical roadmap for future researchers seeking to reproduce or extend this pipeline.

\subsection{Project Overview}
The repository contains the research codebase extending the work of \textbf{Chen Xinnuo} (2025) on evaluating how well synthetic data preserves the clustering properties of real data. The project systematically compares K-Means and Hierarchical Clustering algorithms across 1,800+ datasets to answer critical questions about synthetic data quality regarding:
\begin{itemize}
  \item \textbf{Detection}: Can clustering algorithms correctly identify the number of clusters ($k$) in synthetic data?
  \item \textbf{Robustness}: Which algorithm (K-Means vs. Hierarchical) is more resilient to synthetic data distortion?
  \item \textbf{Quality}: How much does the \texttt{synthpop} CART method degrade cluster geometry and separability?
\end{itemize}

\subsection{Directory Structure}
The repository is logically divided into the original foundational work and the current, automated pipeline:

\begin{verbatim}
ML_SyntheticDataResearch/
|
|-- ofarres/                    <- MAIN RESEARCH PIPELINE (Current Work)
|   |-- SETUP.md                <- Installation & setup instructions
|   |-- README.md               <- Complete pipeline documentation
|   |-- run_all.sh              <- Master execution script
|   |-- requirements.txt        <- Python dependencies
|   |
|   |-- config/                 <- Declarative JSON simulation parameters
|   |-- installer_scripts/      <- Master installers (R + Python packages)
|   |
|   |-- 01_data_generation/     <- Step 1: Generate real + synthetic data
|   |-- 02_clustering/          <- Step 2: Run clustering simulations
|   \-- 03_analysis/            <- Step 3: Visualization & results extraction
|
|-- xinnuo_files/               <- ORIGINAL RESEARCH (Chen Xinnuo, 2025)
\-- README.md                   <- Repository entry point
\end{verbatim}

\subsection{Execution Pipeline}

The environment relies on \textbf{Conda}, \textbf{R (4.0+)}, and \textbf{Python (3.9+)}. The entire simulation environment can be bootstrapped and executed autonomously:

\vspace{0.3cm}
\noindent\textbf{1. Environment Setup}
\begin{verbatim}
conda create -n synthetic_data python=3.9
conda activate synthetic_data
cd ofarres/
./installer_scripts/install_all.sh
\end{verbatim}

\vspace{0.3cm}
\noindent\textbf{2. Pipeline Execution}
\begin{verbatim}
./run_all.sh
\end{verbatim}

\subsection{Pipeline Stages and Outputs}
The automated pipeline executes sequentially through three major stages:

\begin{enumerate}
    \item \textbf{Step 1: Data Generation}
    \begin{itemize}
        \item Generates 180 real datasets with controlled cluster properties (varying $N$, $k$, separation, correlation).
        \item Creates 10 synthetic versions of each real dataset using the \texttt{synthpop} CART method.
        \item \textbf{Output}: 1,980 total datasets (180 real + 1,800 synthetic).
    \end{itemize}
    \item \textbf{Step 2: Clustering Simulation}
    \begin{itemize}
        \item Tests both K-Means and Hierarchical Clustering on all datasets.
        \item Evaluates cluster detection accuracy and quality metrics (e.g., silhouette, distortion).
        \item \textbf{Output}: \texttt{clustering\_results\_final.csv}.
    \end{itemize}
    \item \textbf{Step 3: Visualization \& Analysis}
    \begin{itemize}
        \item Generates publication-ready plots for Success Rate, Performance Delta, and Quality Metrics.
        \item \textbf{Output}: 24 high-resolution PNG visualizations located in \texttt{03\_analysis/}.
    \end{itemize}
\end{enumerate}

\subsection{Citation and Attribution}
This work directly extends the foundational methodology developed by:
\begin{quote}
    \textbf{Chen Xinnuo} (2025). \textit{TFG-EST: Synthetic Data Generation and Clustering Analysis}. Universidad de Barcelona.
\end{quote}
The current extension features automated 3-step pipeline implementation across both R and Python, systematic parameter sweeping, and a reproducible, version-controlled research infrastructure.